\IfFileExists{stacks-project.cls}{%
\documentclass{stacks-project}
}{%
\documentclass{amsart}
}

% For dealing with references we use the comment environment
\usepackage{verbatim}
\newenvironment{reference}{\comment}{\endcomment}
%\newenvironment{reference}{}{}
\newenvironment{slogan}{\comment}{\endcomment}
\newenvironment{history}{\comment}{\endcomment}

% For commutative diagrams we use Xy-pic
\usepackage[all]{xy}

% We use 2cell for 2-commutative diagrams.
\xyoption{2cell}
\UseAllTwocells

% We use multicol for the list of chapters between chapters
\usepackage{multicol}

% This is generally recommended for better output
\usepackage{lmodern}
\usepackage[T1]{fontenc}

% For cross-file-references
\usepackage{xr-hyper}

% Package for hypertext links:
\usepackage{hyperref}

% For any local file, say "hello.tex" you want to link to please
% use \externaldocument[hello-]{hello}
\externaldocument[introduction-]{introduction}
\externaldocument[conventions-]{conventions}
\externaldocument[sets-]{sets}
\externaldocument[categories-]{categories}
\externaldocument[topology-]{topology}
\externaldocument[sheaves-]{sheaves}
\externaldocument[sites-]{sites}
\externaldocument[stacks-]{stacks}
\externaldocument[fields-]{fields}
\externaldocument[algebra-]{algebra}
\externaldocument[brauer-]{brauer}
\externaldocument[homology-]{homology}
\externaldocument[derived-]{derived}
\externaldocument[simplicial-]{simplicial}
\externaldocument[more-algebra-]{more-algebra}
\externaldocument[smoothing-]{smoothing}
\externaldocument[modules-]{modules}
\externaldocument[sites-modules-]{sites-modules}
\externaldocument[injectives-]{injectives}
\externaldocument[cohomology-]{cohomology}
\externaldocument[sites-cohomology-]{sites-cohomology}
\externaldocument[dga-]{dga}
\externaldocument[dpa-]{dpa}
\externaldocument[sdga-]{sdga}
\externaldocument[hypercovering-]{hypercovering}
\externaldocument[schemes-]{schemes}
\externaldocument[constructions-]{constructions}
\externaldocument[properties-]{properties}
\externaldocument[morphisms-]{morphisms}
\externaldocument[coherent-]{coherent}
\externaldocument[divisors-]{divisors}
\externaldocument[limits-]{limits}
\externaldocument[varieties-]{varieties}
\externaldocument[topologies-]{topologies}
\externaldocument[descent-]{descent}
\externaldocument[perfect-]{perfect}
\externaldocument[more-morphisms-]{more-morphisms}
\externaldocument[flat-]{flat}
\externaldocument[groupoids-]{groupoids}
\externaldocument[more-groupoids-]{more-groupoids}
\externaldocument[etale-]{etale}
\externaldocument[chow-]{chow}
\externaldocument[intersection-]{intersection}
\externaldocument[pic-]{pic}
\externaldocument[weil-]{weil}
\externaldocument[adequate-]{adequate}
\externaldocument[dualizing-]{dualizing}
\externaldocument[duality-]{duality}
\externaldocument[discriminant-]{discriminant}
\externaldocument[derham-]{derham}
\externaldocument[local-cohomology-]{local-cohomology}
\externaldocument[algebraization-]{algebraization}
\externaldocument[curves-]{curves}
\externaldocument[resolve-]{resolve}
\externaldocument[models-]{models}
\externaldocument[functors-]{functors}
\externaldocument[equiv-]{equiv}
\externaldocument[pione-]{pione}
\externaldocument[etale-cohomology-]{etale-cohomology}
\externaldocument[proetale-]{proetale}
\externaldocument[relative-cycles-]{relative-cycles}
\externaldocument[more-etale-]{more-etale}
\externaldocument[trace-]{trace}
\externaldocument[crystalline-]{crystalline}
\externaldocument[spaces-]{spaces}
\externaldocument[spaces-properties-]{spaces-properties}
\externaldocument[spaces-morphisms-]{spaces-morphisms}
\externaldocument[decent-spaces-]{decent-spaces}
\externaldocument[spaces-cohomology-]{spaces-cohomology}
\externaldocument[spaces-limits-]{spaces-limits}
\externaldocument[spaces-divisors-]{spaces-divisors}
\externaldocument[spaces-over-fields-]{spaces-over-fields}
\externaldocument[spaces-topologies-]{spaces-topologies}
\externaldocument[spaces-descent-]{spaces-descent}
\externaldocument[spaces-perfect-]{spaces-perfect}
\externaldocument[spaces-more-morphisms-]{spaces-more-morphisms}
\externaldocument[spaces-flat-]{spaces-flat}
\externaldocument[spaces-groupoids-]{spaces-groupoids}
\externaldocument[spaces-more-groupoids-]{spaces-more-groupoids}
\externaldocument[bootstrap-]{bootstrap}
\externaldocument[spaces-pushouts-]{spaces-pushouts}
\externaldocument[spaces-chow-]{spaces-chow}
\externaldocument[groupoids-quotients-]{groupoids-quotients}
\externaldocument[spaces-more-cohomology-]{spaces-more-cohomology}
\externaldocument[spaces-simplicial-]{spaces-simplicial}
\externaldocument[spaces-duality-]{spaces-duality}
\externaldocument[formal-spaces-]{formal-spaces}
\externaldocument[restricted-]{restricted}
\externaldocument[spaces-resolve-]{spaces-resolve}
\externaldocument[formal-defos-]{formal-defos}
\externaldocument[defos-]{defos}
\externaldocument[cotangent-]{cotangent}
\externaldocument[examples-defos-]{examples-defos}
\externaldocument[algebraic-]{algebraic}
\externaldocument[examples-stacks-]{examples-stacks}
\externaldocument[stacks-sheaves-]{stacks-sheaves}
\externaldocument[criteria-]{criteria}
\externaldocument[artin-]{artin}
\externaldocument[quot-]{quot}
\externaldocument[stacks-properties-]{stacks-properties}
\externaldocument[stacks-morphisms-]{stacks-morphisms}
\externaldocument[stacks-limits-]{stacks-limits}
\externaldocument[stacks-cohomology-]{stacks-cohomology}
\externaldocument[stacks-perfect-]{stacks-perfect}
\externaldocument[stacks-introduction-]{stacks-introduction}
\externaldocument[stacks-more-morphisms-]{stacks-more-morphisms}
\externaldocument[stacks-geometry-]{stacks-geometry}
\externaldocument[moduli-]{moduli}
\externaldocument[moduli-curves-]{moduli-curves}
\externaldocument[examples-]{examples}
\externaldocument[exercises-]{exercises}
\externaldocument[guide-]{guide}
\externaldocument[desirables-]{desirables}
\externaldocument[coding-]{coding}
\externaldocument[obsolete-]{obsolete}
\externaldocument[fdl-]{fdl}
\externaldocument[index-]{index}

% Theorem environments.
%
\theoremstyle{plain}
\newtheorem{theorem}[subsection]{Theorem}
\newtheorem{proposition}[subsection]{Proposition}
\newtheorem{lemma}[subsection]{Lemma}

\theoremstyle{definition}
\newtheorem{definition}[subsection]{Definition}
\newtheorem{example}[subsection]{Example}
\newtheorem{exercise}[subsection]{Exercise}
\newtheorem{situation}[subsection]{Situation}

\theoremstyle{remark}
\newtheorem{remark}[subsection]{Remark}
\newtheorem{remarks}[subsection]{Remarks}

\numberwithin{equation}{subsection}

% Macros
%
\def\lim{\mathop{\mathrm{lim}}\nolimits}
\def\colim{\mathop{\mathrm{colim}}\nolimits}
\def\Spec{\mathop{\mathrm{Spec}}}
\def\Hom{\mathop{\mathrm{Hom}}\nolimits}
\def\Ext{\mathop{\mathrm{Ext}}\nolimits}
\def\SheafHom{\mathop{\mathcal{H}\!\mathit{om}}\nolimits}
\def\SheafExt{\mathop{\mathcal{E}\!\mathit{xt}}\nolimits}
\def\Sch{\mathit{Sch}}
\def\Mor{\mathop{\mathrm{Mor}}\nolimits}
\def\Ob{\mathop{\mathrm{Ob}}\nolimits}
\def\Sh{\mathop{\mathit{Sh}}\nolimits}
\def\NL{\mathop{N\!L}\nolimits}
\def\CH{\mathop{\mathrm{CH}}\nolimits}
\def\proetale{{pro\text{-}\acute{e}tale}}
\def\etale{{\acute{e}tale}}
\def\QCoh{\mathit{QCoh}}
\def\Ker{\mathop{\mathrm{Ker}}}
\def\Im{\mathop{\mathrm{Im}}}
\def\Coker{\mathop{\mathrm{Coker}}}
\def\Coim{\mathop{\mathrm{Coim}}}

% Boxtimes
%
\DeclareMathSymbol{\boxtimes}{\mathbin}{AMSa}{"02}

%
% Macros for moduli stacks/spaces
%
\def\QCohstack{\mathcal{QC}\!\mathit{oh}}
\def\Cohstack{\mathcal{C}\!\mathit{oh}}
\def\Spacesstack{\mathcal{S}\!\mathit{paces}}
\def\Quotfunctor{\mathrm{Quot}}
\def\Hilbfunctor{\mathrm{Hilb}}
\def\Curvesstack{\mathcal{C}\!\mathit{urves}}
\def\Polarizedstack{\mathcal{P}\!\mathit{olarized}}
\def\Complexesstack{\mathcal{C}\!\mathit{omplexes}}
% \Pic is the operator that assigns to X its picard group, usage \Pic(X)
% \Picardstack_{X/B} denotes the Picard stack of X over B
% \Picardfunctor_{X/B} denotes the Picard functor of X over B
\def\Pic{\mathop{\mathrm{Pic}}\nolimits}
\def\Picardstack{\mathcal{P}\!\mathit{ic}}
\def\Picardfunctor{\mathrm{Pic}}
\def\Deformationcategory{\mathcal{D}\!\mathit{ef}}
\title{Functors preserve homotopic morphisms of simplicial and cosimplicial objects}
\author{The Stacks Project Authors}
\date{stacks-project:simplicial.tex:local-snapshot-20261012; GNU FDL 1.2 or later}
\begin{document}
\maketitle
\noindent\textit{Editorial note.} The counted statement is the author's own lemma environment labelled lemma-functorial-homotopy (source lines 5222-5249); the proof environment begins on the second line after it and is its own proof, closing at source line 5256.\par\medskip
\noindent\textit{Editorial note.} The proof is a two-sentence pointer: it reduces to covariant functors by the comparison lemma and refers to the author's remark for the functoriality of the homotopy relation. The card reproduces that as written and carries the remark and the comparison lemma as context, so nothing is silently completed.\par\medskip
\noindent\textit{Editorial note.} The retained preamble is the chapter's own minimal preamble (sources/stacks-simplicial/preamble.tex, taken with its real line count), a verbatim excerpt of the native Xy-pic setup, the native theorem declarations with the mathematical macros, and the project's own no-op reference environments; no macro was invented or redefined.\par\medskip
\noindent\textit{Editorial scope.} One original result of the Stacks Project's Simplicial Methods chapter, counted once: any functor preserves homotopic morphisms of simplicial and cosimplicial objects, in all four variance cases.\par\medskip
\noindent\textit{Allowed standard prerequisites (editorial).}\begin{itemize}
\item Simplicial and cosimplicial objects. The categories of simplicial and of cosimplicial objects in a category, with their morphisms and the notion of homotopy between them.. They are the objects of the four clauses of the statement (source lines 5224-5247).
\item Homotopy in terms of the comparison categories. The description of the homotopy relation through the auxiliary categories of the author's remarks, which makes the relation manifestly functorial (the author's remark-homotopy-better and remark-homotopy-cosimplicial-better, retained as context).. It is where the substance of the proof lives (source lines 5254-5255).
\item Reduction of a contravariant functor to a covariant one. A contravariant functor may be regarded as a covariant functor between the appropriate pair of categories (the author's lemma-compare-homotopies, retained as context).. It reduces the four cases to the covariant ones (source lines 5252-5253).
\end{itemize}\medskip
\begin{remark}
\label{remark-homotopy-cosimplicial-better}
Let $\mathcal{C}$ be any category (no assumptions whatsoever). Let
$U$ and $V$ be cosimplicial objects of $\mathcal{C}$. Let $a, b : U \to V$
be morphisms of cosimplicial objects of $\mathcal{C}$. A
{\it homotopy from $a$ to $b$} is given by morphisms
$h_{n, \alpha} : U_n \to V_n$, for $n \geq 0$, $\alpha \in \Delta[1]_n$
satisfying (\ref{equation-property-homotopy-cosimplicial})
for all morphisms $f$ of $\Delta$ and such that
$a_n = h_{n, 0 : [n] \to [1]}$ and $b_n = h_{n, 1 : [n] \to [1]}$
for all $n \geq 0$.
As in Definition \ref{definition-homotopy-cosimplicial}
we say the morphisms $a$ and $b$
are {\it homotopic} if there exists a sequence of morphisms
$a = a_0, a_1, \ldots, a_n = b$ from $U$ to $V$ such that for each
$i = 1, \ldots, n$ there either exists a homotopy from $a_{i - 1}$ to $a_i$
or there exists a homotopy from $a_i$ to $a_{i - 1}$.
Clearly, if $F : \mathcal{C} \to \mathcal{C}'$ is any functor
and $\{h_{n, i}\}$ is a homotopy from $a$ to $b$, then
$\{F(h_{n, i})\}$ is a homotopy from $F(a)$ to $F(b)$.
Similarly, if $a$ and $b$ are homotopic, then $F(a)$ and $F(b)$
are homotopic. This new notion is the same
as the old one in case finite products exist. We deduce
in particular that functors preserve the original notion
whenever both categories have finite products.
\end{remark}

\begin{remark}
\label{remark-homotopy-better}
Let $\mathcal{C}$ be any category (no assumptions whatsoever). Let
$U$ and $V$ be simplicial objects of $\mathcal{C}$. Let $a, b : U \to V$
be morphisms of simplicial objects of $\mathcal{C}$. A
{\it homotopy from $a$ to $b$} is given by
morphisms\footnote{In the literature, often the maps
$h_{n + 1, i} \circ s_i : U_n \to V_{n + 1}$ are used instead
of the maps $h_{n, i}$. Of course the relations these maps satisfy
are different from the ones in Lemma \ref{lemma-relations-homotopy}.}
$h_{n, i} : U_n \to V_n$, for $n \geq 0$, $i = 0, \ldots, n + 1$
satisfying the relations of Lemma \ref{lemma-relations-homotopy}.
As in Definition \ref{definition-homotopy} we say the morphisms $a$ and $b$
are {\it homotopic} if there exists a sequence of morphisms
$a = a_0, a_1, \ldots, a_n = b$ from $U$ to $V$ such that for each
$i = 1, \ldots, n$ there either exists a homotopy from $a_{i - 1}$ to $a_i$
or there exists a homotopy from $a_i$ to $a_{i - 1}$.
Clearly, if $F : \mathcal{C} \to \mathcal{C}'$ is any functor
and $\{h_{n, i}\}$ is a homotopy from $a$ to $b$, then
$\{F(h_{n, i})\}$ is a homotopy from $F(a)$ to $F(b)$.
Similarly, if $a$ and $b$ are homotopic, then $F(a)$ and $F(b)$
are homotopic.
Since the lemma says that the newer notion is the same
as the old one in case finite coproduct exist, we deduce
in particular that functors preserve the original notion
whenever both categories have finite coproducts.
\end{remark}

\begin{lemma}
\label{lemma-compare-homotopies}
Let $\mathcal{C}$ be a category. Suppose that $U$ and $V$ are two
cosimplicial objects of $\mathcal{C}$. Let $a, b : U \to V$ be morphisms
of cosimplicial objects. Recall that $U$, $V$ correspond
to simplicial objects $U'$, $V'$ of $\mathcal{C}^{opp}$.
Moreover $a, b$ correspond to morphisms $a', b' : V' \to U'$.
The following are equivalent
\begin{enumerate}
\item There exists a homotopy $h = \{h_{n, \alpha}\}$ from
$a$ to $b$ as in Remark \ref{remark-homotopy-cosimplicial-better}.
\item There exists a homotopy $h = \{h_{n, i}\}$ from $a'$ to $b'$
as in Remark \ref{remark-homotopy-better}.
\end{enumerate}
Thus $a$ is homotopic to $b$ as in
Remark \ref{remark-homotopy-cosimplicial-better}
if and only if $a'$ is homotopic to $b'$ as in
Remark \ref{remark-homotopy-better}.
\end{lemma}

\begin{proof}
In case $\mathcal{C}$ has finite products, then $\mathcal{C}^{opp}$
has finite coproducts and we may use
Definitions \ref{definition-homotopy-cosimplicial}
and \ref{definition-homotopy} instead of
Remarks \ref{remark-homotopy-cosimplicial-better} and
\ref{remark-homotopy-better}.
In this case $h : U \to \Hom(\Delta[1], V)$
is the same as a morphism $h' : \Hom(\Delta[1], V)' \to U'$.
Since products and coproducts get switched too, it is immediate
that $(\Hom(\Delta[1], V))' = V' \times \Delta[1]$.
Moreover, the ``primed'' version of the morphisms
$e_0, e_1 : \Hom(\Delta[1], V) \to V$ are the morphisms
$e_0, e_1 : V' \to \Delta[1] \times V$.
Thus $e_0 \circ h = a$ translates into $h' \circ e_0 = a'$ and
similarly $e_1 \circ h = b$ translates into $h' \circ e_1 = b'$.
This proves the lemma in this case.

\medskip\noindent
In the general case, one needs to translate the relations given
by (\ref{equation-property-homotopy-cosimplicial}) into the relations
given in Lemma \ref{lemma-relations-homotopy}. We omit the details.

\medskip\noindent
The final assertion is formal from the equivalence of (1) and (2).
\end{proof}

\begin{definition}
\label{definition-homotopy-cosimplicial}
Let $\mathcal{C}$ be a category having finite products.
Let $U$ and $V$ be two cosimplicial objects of $\mathcal{C}$.
Let $a, b : U \to V$ be two morphisms of cosimplicial objects
of $\mathcal{C}$.
\begin{enumerate}
\item We say a morphism
$$
h : U \longrightarrow \Hom(\Delta[1], V)
$$
such that $a = e_0 \circ h$ and $b = e_1 \circ h$ is a
{\it homotopy from $a$ to $b$}.
\item We say $a$ and $b$ are {\it homotopic} or are
{\it in the same homotopy class} if there exists a sequence
$a = a_0, a_1, \ldots, a_n = b$ of morphisms from $U$ to $V$
such that for each $i = 1, \ldots, n$ there either exists a
homotopy from $a_i$ to $a_{i - 1}$ or there exists a homotopy
from $a_{i - 1}$ to $a_i$.
\end{enumerate}
\end{definition}

\begin{equation}
\label{equation-property-homotopy-cosimplicial}
h_{m, \alpha} \circ U(f) = V(f) \circ h_{n, \alpha \circ f}
\end{equation}

\begin{definition}
\label{definition-homotopy}
Let $\mathcal{C}$ be a category having finite coproducts.
Suppose that $U$ and $V$ are two simplicial objects of $\mathcal{C}$.
Let $a, b : U \to V$ be two morphisms.
\begin{enumerate}
\item We say a morphism
$$
h : U \times \Delta[1] \longrightarrow V
$$
is a {\it homotopy from $a$ to $b$} if $a = h \circ e_0$ and
$b = h \circ e_1$.
\item We say the morphisms $a$ and $b$ are {\it homotopic} or are
{\it in the same homotopy class}
if there exists a sequence of morphisms $a = a_0, a_1, \ldots, a_n = b$
from $U$ to $V$ such that for each $i = 1, \ldots, n$ there either exists
a homotopy from $a_{i - 1}$ to $a_i$ or there exists a homotopy
from $a_i$ to $a_{i - 1}$.
\end{enumerate}
\end{definition}

\begin{lemma}
\label{lemma-relations-homotopy}
In the situation above, we have the following relations:
\begin{enumerate}
\item We have $h_{n, 0} = b_n$ and $h_{n, n + 1} = a_n$.
\item We have $d^n_j \circ h_{n, i} = h_{n - 1, i - 1} \circ d^n_j$
for $i > j$.
\item We have $d^n_j \circ h_{n, i} = h_{n - 1, i} \circ d^n_j$
for $i \leq j$.
\item We have $s^n_j \circ h_{n, i} = h_{n + 1, i + 1} \circ s^n_j$
for $i > j$.
\item We have $s^n_j \circ h_{n, i} = h_{n + 1, i} \circ s^n_j$
for $i \leq j$.
\end{enumerate}
Conversely, given a system of maps $h_{n, i}$ satisfying the
properties listed above, then these define a morphism
$h$ which is a homotopy from $a$ to $b$.
\end{lemma}

\begin{proof}
Omitted. You can prove the last statement using the fact,
see Lemma \ref{lemma-face-degeneracy-category} that
to give a morphism of simplicial objects is the
same as giving a sequence of morphisms $h_n$ commuting
with all $d^n_j$ and $s^n_j$.
\end{proof}

\begin{lemma}
\label{lemma-face-degeneracy-category}
The category $\Delta$ is the universal category
with objects $[n]$, $n \geq 0$ and morphisms
$\delta^n_j$ and $\sigma^n_j$ such that (a) every morphism is
a composition of these morphisms, (b) the relations
listed in Lemma \ref{lemma-relations-face-degeneracy} are satisfied,
and (c) any relation among the morphisms is a consequence of
those relations.
\end{lemma}

\begin{proof}
Omitted.
\end{proof}

\begin{lemma}
\label{lemma-relations-face-degeneracy}
The morphisms $\delta^n_j$ and $\sigma^n_j$ satisfy the following relations.
\begin{enumerate}
\item If $0 \leq i < j \leq n + 1$, then
$\delta^{n + 1}_j \circ \delta^n_i =
\delta^{n + 1}_i \circ \delta^n_{j - 1}$.
In other words the diagram
$$
\xymatrix{
& [n] \ar[rd]^{\delta^{n + 1}_j} & \\
[n - 1] \ar[ru]^{\delta^n_i} \ar[rd]_{\delta^n_{j - 1}} & &
[n + 1] \\
& [n] \ar[ru]_{\delta^{n + 1}_i} &
}
$$
commutes.
\item If $0 \leq i < j \leq n - 1$, then
$\sigma^{n - 1}_j \circ \delta^n_i =
\delta^{n - 1}_i \circ \sigma^{n - 2}_{j - 1}$.
In other words the diagram
$$
\xymatrix{
& [n] \ar[rd]^{\sigma^{n - 1}_j} & \\
[n - 1] \ar[ru]^{\delta^n_i} \ar[rd]_{\sigma^{n - 2}_{j - 1}} & &
[n - 1] \\
& [n - 2] \ar[ru]_{\delta^{n - 1}_i} &
}
$$
commutes.
\item If $0 \leq j \leq n - 1$, then
$\sigma^{n - 1}_j \circ \delta^n_j = \text{id}_{[n - 1]}$
and
$\sigma^{n - 1}_j \circ \delta^n_{j + 1} = \text{id}_{[n - 1]}$.
In other words the diagram
$$
\xymatrix{
& [n] \ar[rd]^{\sigma^{n - 1}_j} & \\
[n - 1]
\ar[ru]^{\delta^n_j}
\ar[rd]_{\delta^n_{j + 1}}
\ar[rr]^{\text{id}_{[n - 1]}} & & [n - 1] \\
& [n] \ar[ru]_{\sigma^{n - 1}_j} &
}
$$
commutes.
\item If $0 < j + 1 < i \leq n$, then
$\sigma^{n - 1}_j \circ \delta^n_i =
\delta^{n - 1}_{i - 1} \circ \sigma^{n - 2}_j$.
In other words the diagram
$$
\xymatrix{
& [n] \ar[rd]^{\sigma^{n - 1}_j} & \\
[n - 1] \ar[ru]^{\delta^n_i} \ar[rd]_{\sigma^{n - 2}_j} & &
[n - 1] \\
& [n - 2] \ar[ru]_{\delta^{n - 1}_{i - 1}} &
}
$$
commutes.
\item If $0 \leq i \leq j \leq n - 1$, then
$\sigma^{n - 1}_j \circ \sigma^n_i =
\sigma^{n - 1}_i \circ \sigma^n_{j + 1}$.
In other words the diagram
$$
\xymatrix{
& [n] \ar[rd]^{\sigma^{n - 1}_j} & \\
[n + 1] \ar[ru]^{\sigma^n_i} \ar[rd]_{\sigma^n_{j + 1}} & &
[n - 1] \\
& [n] \ar[ru]_{\sigma^{n - 1}_i} &
}
$$
commutes.
\end{enumerate}
\end{lemma}

\begin{proof}
Omitted.
\end{proof}

\begin{lemma}
\label{lemma-functorial-homotopy}
\begin{slogan}
Functors preserve homotopic morphisms of (co)simplicial objects.
\end{slogan}
Let $\mathcal{C}, \mathcal{C}', \mathcal{D}, \mathcal{D}'$ be categories.
With terminology as in Remarks \ref{remark-homotopy-cosimplicial-better} and
\ref{remark-homotopy-better}.
\begin{enumerate}
\item Let $a, b : U \to V$ be morphisms of simplicial objects
of $\mathcal{D}$. Let $F : \mathcal{D} \to \mathcal{D}'$ be a covariant
functor. If $a$ and $b$ are homotopic, then $F(a)$, $F(b)$
are homotopic morphisms $F(U) \to F(V)$ of simplicial objects.
\item Let $a, b : U \to V$ be morphisms of cosimplicial objects
of $\mathcal{C}$. Let $F : \mathcal{C} \to \mathcal{C}'$ be a covariant
functor. If $a$ and $b$ are homotopic, then $F(a)$, $F(b)$
are homotopic morphisms $F(U) \to F(V)$ of cosimplicial objects.
\item Let $a, b : U \to V$ be morphisms of simplicial objects of $\mathcal{D}$.
Let $F : \mathcal{D} \to \mathcal{C}$ be a contravariant
functor. If $a$ and $b$ are homotopic, then $F(a)$, $F(b)$
are homotopic morphisms $F(V) \to F(U)$ of cosimplicial objects.
\item Let $a, b : U \to V$ be morphisms of cosimplicial objects of
$\mathcal{C}$.
Let $F : \mathcal{C} \to \mathcal{D}$ be a contravariant
functor. If $a$ and $b$ are homotopic, then $F(a)$, $F(b)$
are homotopic morphisms $F(V) \to F(U)$ of simplicial objects.
\end{enumerate}
\end{lemma}

\begin{proof}
By Lemma \ref{lemma-compare-homotopies} above, we can
turn $F$ into a covariant functor between a pair of
categories, and we have to show that the functor preserves homotopic pairs
of maps. This is explained in Remark \ref{remark-homotopy-better}.
\end{proof}

\end{document}
